\chapter{Nodos LoRa}

En términos de redes de telecomunicaciones móviles, a los dispositivos que consumen la red inalámbrica se los denomina comúnmente como \emph{nodos}, y es así que en \gls{lwan} se hablará del \emph{dispositivo \gls{lora}}, también conocido como \emph{mote}. De igual manera, en el modelo cliente/servidor, el nodo se comunicará con el servidor \gls{lwan}, evitando la comunicación directa con otros dispositivos de la red. El servidor \gls{lwan} podrá enviarle de vuelta al nodo mensajes almacenados en el servidor \gls{lora}, siempre dependiendo de la clase de trabajo en que se encuentre el nodo (sección \ref{secc:clases}). Para efectos prácticos, se hablará del nodo funcionando en clase A, en la que el nodo siempre inicia la comunicación. Las demás clases conservan el mismo modelo cliente/servidor, por lo que no se abordarán en detalle.

La composición del nodo se basa típicamente en una arquitectura de \emph{hardware} de referencia (figura \ref{fig:arqgenhw}), la cual se puede abordar según la funcionalidad de cada uno de los bloques o según la implementación disponible de la arquitectura base mencionada.

\section{Bloques de la arquitectura genérica}

En la arquitectura genérica de la figura \ref{fig:arqgenhw}, cada bloque tiene una función específica dentro del nodo:
\begin{itemize}
    \item La fuente de energía puede ser suministrada mediante un conector o una batería.
    \item La \gls{mcu}, o simplemente el \gls{uc}, maneja todas las funciones del dispositivo e implementa la pila LoRaWAN.
    \item El radio LoRa está compuesto por el transceptor \gls{lora}, el circuito de acople de la antena y la antena en sí misma.
    \item Los periféricos pueden ser sensores como acelerómetros o sensores de temperatura, o entradas y salidas tales como relés o pantallas.
\end{itemize}

\begin{figure}[H]
    \caption{Arquitectura de \emph{hardware} de un dispositivo \gls{lora} genérico}
    \label{fig:arqgenhw}
    \centering
    \includegraphics[width=0.8\textwidth]{arq-gen-hard}
    \propia
\end{figure}

\subsection{Fuente de poder}

En un \gls{uc} generalmente se observa que el consumo de energía es directamente proporcional a la velocidad de procesamiento y a la cantidad de periféricos activos, entre otros factores. 

Esto es crítico en aplicaciones que funcionan la mayor parte del tiempo, como una \gls{wsn} (capítulo \ref{cap:aplicacion}), por lo cual se hace conveniente disminuir el consumo. Una forma de lograrlo es haciendo que el nodo pase la mayor parte del tiempo en un estado ``adormecido'', esto es, un estado en que su actividad de procesamiento es mínima. Dependiendo del tipo de \gls{uc}, de la arquitectura y de los diversos estados de adormecimiento, la corriente requerida suele oscilar entre algunos microamperios y cientos de microamperios. El consumo de energía en trabajo normal es una medida útil y se hace más importante cuanto más tiempo pasa un nodo en periodos de trabajo normal, donde los consumos exigen corrientes que oscilan entre \SI{1}{mA} y \SI{10}{mA} para \gls{uc} pequeños y de bajo consumo, pero pueden variar dependiendo de los tipos de periféricos incorporados, la velocidad del reloj, el tamaño de la memoria y la localización del programa de ejecución.

La tarea primordial del nodo sensor es detectar eventos, realizar el procesamiento básico de los datos y transmitir la información. Lo anterior implica que en el consumo de energía se destacan tres aspectos: sensado, procesamiento y comunicación. La comunicación incluye transmisión y recepción, aspectos que demandan mayor consumo. \cite{Akyildiz-2002}.

La energía será usualmente un recurso limitado en el nodo sensor, lo que implica el uso de fuentes de energía alternativas tales como celdas solares o generadores eólicos, dependiendo del tipo de ambiente en el cual se halla el módulo sensor.

Otros aspectos relativos al suministro de energía están ligados al tiempo de vida de las fuentes y al acceso al nodo sensor. 

\subsection{Microcontrolador}

El \gls{uc} se puede considerar como el centro de control del nodo sensor. Casi todas las familias de \gls{uc} disponibles son diferentes: algunos pueden ser diseñados para la velocidad, y otros para la eficiencia energética. En la cuestión de elegir el \gls{uc} correcto para la aplicación particular del nodo, hay varios aspectos para tener en cuenta en el momento de diseñarlo.

A continuación, se presentan los componentes más relevantes del \gls{uc}:

\begin{itemize}
	\item Memoria:    
    la cantidad y el tipo de memoria utilizados son dos variables que están ligadas al consumo de energía y que deben ser cuidadosamente elegidas, dado que ambos parámetros implican costos por controlar. En general, a mayor cantidad de RAM, más potencia requerida para su sostenimiento; no obstante, los \gls{uc} modernos pueden incorporar varios cientos de kilobytes de memoria \emph{flash} y decenas de kilobytes de RAM y continúan utilizando paquetes estandarizados de manejo con costos bajos.
    
    Una tecnología nueva en este aspecto es el uso de FRAM, una alternativa de memoria que combina la velocidad de la memoria RAM convencional con la no volatilidad del almacenamiento de la memoria \emph{flash} y que elimina la separación entre las dos. Este tipo de memoria tiene menor consumo, cuando funciona, en comparación con la memoria \emph{flash} estándar, y elimina la necesidad de refrescamiento de la RAM para preservar su contenido.

    \item Soporte de entrada-salida: 
    el \gls{uc} está en capacidad de llevar a cabo funciones de comunicación tipo serial o paralela, temporizaciones, comparaciones, transmisión y recepción digital y analógica en diversos protocolos y velocidades, además del procesamiento de datos. Una de las limitaciones a estas funciones está en el número de pines de acceso físico a su estructura. Para dar solución a esta situación, se trata de mantener al máximo el número de funciones y se multiplexa la comunicación de  entrada-salida en los pines disponibles, seleccionando a través de \emph{software} la gran variedad de operaciones que permita el dispositivo.

    Para el diseño de un nodo sensor se requieren muchas de las funciones presentes en el \gls{uc} y una cierta cantidad de pines de comunicación con funciones particulares. Como ejemplos de funciones y pines requeridos, se tienen las interfaces \IC, \IS, canales SPI, conversores ADC y DAC, pines de interrupciones E/S, señales de tiempo y de reloj, bits de bandera, entre otros.

    \item Soporte de comunicación: 
    los nodos sensores están normalmente dispersos en áreas específicas de medición. Cada uno de estos nodos tiene la capacidad de recopilar datos y enviarlos a la pasarela y/o a los usuarios finales.

    Los protocolos utilizados por el receptor y por los nodos de sensores permiten el encaminamiento e integración de los datos con los protocolos de red de manera eficiente a través del medio inalámbrico, distribuyendo el trabajo de los nodos sensores.

	Los nodos sensores están debidamente coordinados y administrados para que realicen tareas que disminuyan el consumo total de energía; por ejemplo, informar a sus vecinos cuando su nivel de energía es bajo y por tanto no puede participar en el intercambio de mensajes.

    \item Adquisición de datos: 
    el modelo que se siga para la toma de datos caracteriza y describe la interacción entre los sensores y la aplicación.

    El modelo de adquisición de datos más apropiado para una aplicación se determina según las características que se desee dar a la aplicación; en las diferentes situaciones, los datos fluyen principalmente desde el nodo del sensor hasta la puerta de enlace; en las aplicaciones de control, los datos, adicionalmente fluyen desde la puerta de enlace hacia los nodos sensores.

	Las técnicas de monitoreo más usuales \cite{Cecilio-2014} son:
	\begin{itemize}
 		\item Muestreo periódico: para aplicaciones donde las condiciones o los procesos requieren ser monitoreados constantemente, como es el caso de la temperatura en un ambiente acondicionado o la presión en una tubería de proceso. Para este caso, los datos se adquieren en el nodo sensor y son reenviados a la puerta de enlace de forma periódica. El periodo de muestreo depende de qué tan rápido varía la magnitud monitoreada y/o qué características de control se requieren en el proceso. Si la dinámica del proceso controlado varía, los nodos sensores deben ser capaces de adaptar sus tasas de muestreo a la dinámica del proceso. Esta adaptación debe hacerse automáticamente o a través de comandos emitidos por la aplicación.
		\item Muestreo impuesto por el evento: hay muchos casos que requieren el monitoreo de variables críticas, como las alarmas de fuego, las alarmas de presión o los sensores de puertas o ventanas, etc. La detección de dichos eventos exige que el nodo sensor actúe con eficiencia en el consumo de energía y con la adecuada velocidad de respuesta.
		\item Captura-proceso-almacenamiento: en algunos casos, los datos pueden ser capturados y almacenados o incluso procesados por el nodo sensor antes de que se transmita a la puerta de enlace, en lugar de transmitir inmediatamente cada muestra a medida que se adquiere, lo cual mejora el rendimiento de la red en cuanto a eficiencia de consumo y ancho de banda.
	\end{itemize}
\end{itemize}

\subsection{Radio LoRa}

La comunicación de radio puede ser bidireccional, tanto entre los nodos como hacia las estaciones base. Adicionalmente, se debe buscar un balance entre el consumo de energía del módem \gls{lora} y la \gls{snr}. Dependiendo de este balance, habrá un impacto consecuente en el consumo de energía de este elemento del sistema, comparado con los otros, sobre todo en el momento en que el nodo requiera transmitir información.

\subsection{Periféricos}

Los sensores y actuadores son la razón de ser del nodo. Estos se distribuyen dentro del área de cobertura del nodo para supervisar y recopilar datos con el objetivo de transmitirlos a otros nodos y/o a un servidor.

\section{Topologías}

Dependiendo de los requisitos de diseño y producción, se tienen varias opciones para construir un dispositivo \gls{lora}. La elección de la arquitectura objetivo se hace con base en criterios como las cantidades de producción esperadas, las habilidades disponibles en ingeniería de RF y la línea de tiempo de desarrollo disponible para completar el proyecto.

Con el fin de seleccionar la arquitectura óptima, es preferible partir de un kit de desarrollo \gls{lora} y crear un bosquejo.

Por otro lado, en la figura \ref{fig:arqgensw} se muestra una arquitectura genérica del \emph{software} en el dispositivo.

\begin{figure}[t]
    \centering
    \caption{Arquitectura genérica del \emph{software} del nodo}
    \label{fig:arqgensw}
    \includegraphics[width=\textwidth]{arq-gen-soft}
    \propia
\end{figure}

La capa de \emph{drivers} provee la adaptación de \emph{hardware} e implementa todos los \emph{drivers} y \emph{buffers} para controlar los dispositivos periféricos. Esta permite abstraer el \emph{hardware} como funciones simples expuestas al \emph{middleware}.

El \emph{middleware} implementa las librerías de protocolos de comunicación (\gls{lwan}, 6LowPAN, etc.) y los controladores complejos, como el control de pantallas y de GPS.

La capa de aplicación contiene todas las aplicaciones funcionales en las que se implementa el comportamiento del dispositivo y sus funcionalidades.

\subsection{Diseño basado en un chipset LoRa}

Este diseño sigue la arquitectura genérica de la figura \ref{fig:arqgenhw}, con el radio \gls{lora} basado en un transceptor \gls{lora} de Semtech y la antena, más su circuitería asociada.

\begin{figure}[t]
    \centering
    \caption{Arquitectura de dispositivo basado en un chipset LoRa}
    \label{fig:loraWchipset}
    \includegraphics[width=\textwidth]{design-chipset}
    \propia
\end{figure}

En esta arquitectura, ilustrada en la figura \ref{fig:loraWchipset}, el desarrollador del dispositivo deberá tener cuidado de todo el diseño RF, incluyendo las restricciones de enrutamiento en la PCB, el acoplamiento y sintonización de la antena y los problemas de emisión e inmunidad.

La totalidad de la pila \gls{lwan} será llevada por el MCU principal e implementada por el desarrollador o fabricante del dispositivo.

\subsection{Diseño basado en un módulo LoRa}

El módulo \gls{lora} es un componente que contiene el \gls{mcu} y el radio \gls{lora}. Se tiene la disponibilidad de programar el \emph{software} de su \gls{mcu} para ejecutar la aplicación y la pila \gls{lwan}.

El principal beneficio de este enfoque de diseño es que todo el desarrollo del \emph{hardware} de RF es implementado por el fabricante del módulo.

La mayor parte de la sintonización y acoplamiento de la antena es realizada dentro del módulo. El fabricante del módulo brinda un diseño de referencia para conectar o rediseñar la antena.

\begin{figure}[t]
    \centering
    \caption{Arquitectura de dispositivo basado en un módulo LoRa}
    \label{fig:loraWmodule}
    \includegraphics[width=\textwidth]{design-module}
    \propia
\end{figure}

En comparación con la arquitectura anterior, dado que se dispone del MCU, en esta arquitectura (figura~\ref{fig:loraWmodule}), el fabricante del dispositivo se encarga solamente de la implementación de la pila \gls{lwan}. Dependiendo del proveedor del módulo, la pila \gls{lwan} se puede entregar como una librería.

\subsection{Diseño basado en un RF-MCU}

Un RF-MCU es un \gls{soc} que incluye un \gls{mcu} y un transceptor LoRa en el mismo circuito integrado. Los requerimientos RF son los mismos del diseño RF anterior, pero esta arquitectura ofrece mayor ganancia en el tamaño total del dispositivo. Desde el punto de vista del desarrollo, este diseño es similar al del módulo.

\subsection{Diseño basado en un módem LoRa}

Un módem LoRa es un componente que contiene todos los elementos de radio: pila y circuitería RF (figura \ref{fig:loraModemArch}). Algunos módems pueden también incluir una antena integrada, o el fabricante da una referencia de diseño para diseñar o conectar una antena externa.

\begin{figure}[t]
    \centering
    \caption{Arquitectura de dispositivo basado en módem LoRa}
    \label{fig:loraModemArch}
    \includegraphics[width=0.9\textwidth]{design-modem}
    \propia
\end{figure}

El módem integra la pila \gls{lwan} entera y actúa como esclavo, de tal manera que requiere de un \emph{host} para manejarlo. La comunicación usualmente se realiza mediante comandos AT para configurar o enviar mensajes.

La interfaz de \emph{hardware} puede ser usualmente UART, USB o SPI.

\subsection{Dispositivos con un módem LoRa externo}

Un módem LoRa también puede ser un dispositivo independiente externo (figura~\ref{fig:loraExternal}), conectado a un dispositivo existente mediante una conexión USB, por ejemplo.

En ese caso, el dispositivo existente debe tener una conexión externa disponible y también debe ser capaz de energizar el módem LoRa externo.

El módem integra la pila \gls{lwan} entera y actúa como un esclavo, de tal modo que requiere de un controlador externo para manejarlo. La comunicación usualmente se realiza mediante comandos AT para configurarlo o enviar mensajes.

\begin{figure}[H]
    \centering
    \caption{Arquitectura de dispositivo basado en módem LoRa externo}
    \label{fig:loraExternal}
    \includegraphics[width=0.9\textwidth]{design-external}
    \propia
\end{figure}

\section{Interfaces digitales}

Los sensores disponibles en el mercado pueden requerir polarización y acondicionamiento, como pueden ya tener integrados estos circuitos en el mismo encapsulado e inclusive presentar una salida con alguna interfaz digital comercial.

Debido a la variedad tanto de interfaces digitales comerciales como de distintas configuraciones de polarización para los sensores que requieren circuitos de polarización y/o acondicionamiento, se plantea una interfaz que permita conectar la mayor cantidad posible de componentes necesarios entre el \gls{uc} conectado al nodo LoRaWAN y el sensor que se vaya a conectar al nodo.

Los protocolos de comunicación electrónica más comunes entre componentes de sistemas embebidos suelen ser mediante buses serie de varios tipos. Se hablará solamente de las líneas dedicadas exclusivamente al protocolo en mención, excluyendo, así, las de polarización y energización de los dispositivos.

\subsection{Inter-Integrated Circuits (I2C)}

El bus \gls{i2c} se desarrolló en 1982 por Phillips (actualmente, NXP). Se utiliza principalmente en la comunicación de dispositivos distribuidos en un mismo sistema embebido, por ejemplo, entre un \gls{uc} y circuitos periféricos integrados. Es un bus maestro-esclavo, en el que la transferencia de datos es siempre inicializada por un maestro y el esclavo reacciona ante esta (figura~\ref{fig:i2c}).

\begin{quote}\small
    Es posible tener varios maestros mediante un modo multimaestro, en el que se pueden comunicar dos maestros entre sí, de modo que uno de ellos trabaja como esclavo. El arbitraje (control de acceso al bus) se rige por las especificaciones, de modo que los maestros pueden ir turnándose. (Aguilar Castillejos, 2019, p. 1)
\end{quote}
    
Requiere solamente dos líneas, para reloj y datos. Ambas se encuentran en un estado lógico alto débil (debido a resistencias \emph{pull up} en cada línea), y por lo mismo quienes usen el bus han de tener salidas tipo colector/drenador abierto. Quien inicia una transmisión (tomando el control de las dos líneas) es siempre un dispositivo maestro, y puede iniciarla con uno de los demás dispositivos en el bus, que se comportarán como esclavos dependiendo del direccionamiento.

\begin{figure}[H]
    \centering
    \caption{Esquema \IC}\label{fig:i2c}
    
    \includegraphics[width=\textwidth]{I2C}

    \nota{realmente, la polarización (líneas VCC y GND) puede ser distinta para cada dispositivo, pero las líneas del bus, SDA y SCL, son comunes a todos eléctricamente.}
    
    \vspace{.3\baselineskip}
    \fuente{\url{http://www.prometec.net/bus-i2c/}}
\end{figure}

\subsection{SPI}

También en modelo maestro-esclavo, con comunicación iniciada por el dispositivo maestro, este bus permite comunicación \emph{full duplex} mediante las líneas MOSI (\emph{master out, slave in}) y MISO (\emph{master in, slave out}), con reloj y habilitación del esclavo establecidos por el maestro, mediante las líneas CLK y SS (\emph{slave select}) (figura~\ref{fig:spi}).

Este bus permite velocidades de transmisión más rápidas, tanto porque las líneas son activas (requieren salidas tipo \emph{push-pull}) como porque el maestro puede enviar y recibir al tiempo, a costa de requerir más pines que \IC\ o UART.

\begin{figure}
    \centering
    \caption{Esquema SPI}\label{fig:spi}
    \includegraphics{SPI}
    \nota{todos los esclavos comparten y se conectan a las mismas líneas del bus: SCLK, MOSI y MISO, pero solo son activos bajo su señal SS.}

    \vspace{.3\baselineskip}
    \fuente{Cburnett (2006).}
\end{figure}

\subsection{UART}

El transmisor-receptor asíncrono universal\foreignlanguage{english}{(\emph{universal asynchronous receiver-transmitter})} es similar a SPI dado que permite comunicación \emph{full duplex}, mas no utiliza ni las líneas de reloj (dado que cada bit tiene un tiempo determinado y en mutuo acuerdo) ni de habilitación (cualquiera de los dos puede iniciar una transmisión). Como mínimo, cada dispositivo que soporte UART presenta las líneas Rx y Tx; para comunicar dos dispositivos UART basta con intercambiar las líneas (figura~\ref{fig:uart}).

\begin{figure}[t]
    \centering
    \caption{Esquema UART}
    \label{fig:uart}
    \includegraphics[width=0.6\textwidth]{uart}
    \nota{las líneas de los dispositivos se conectan intercambiadas, emparejando siempre un receptor con un transmisor.}
    
    \vspace{.3\baselineskip}
    \fuente{\url{https://www.mikroe.com/blog/uart-serial-communication}}
\end{figure}

\section{Interfaz planteada}

La interfaz abarca dos entornos: uno de \emph{hardware}, relacionado con la conexión física, y otro de \emph{software}, relacionado con una capa de abstracción entre los distintos métodos de obtención de las variables de los sensores, ya sea mediante protocolos digitales o mediante una cadena de medida, con el fin de permitir su transmisión sobre una aplicación en la red \gls{lwan}.

Con ayuda de una aplicación \gls{lwan}, esta interfaz permitirá la interconexión de dos circuitos: uno de estos será el núcleo Pylates, conformado por un \gls{mcu} y un módem \gls{lwan}. El otro será la placa de aplicación, con todos los circuitos adicionales requeridos en el nodo.

\subsection{Interfaz de \emph{hardware}: conector y funcionalidades de pines}

Ha de ser posible que el \gls{mcu} acceda a las variables de medida de cada sensor, bien mediante interfaces digitales (SPI, \IC, UART), bien con la implementación de una cadena de medida (OPAMPS+REF+ADC+DAC) para aquellos en que sea necesario. Se estima el peor de los escenarios, en el que se requiere la mayor cantidad de pines: cuatro sensores conectados al \gls{mcu}, tres mediante las interfaces digitales mencionadas y uno con una cadena de medida asistida por los periféricos analógicos del \gls{mcu}. Otra alternativa consiste en considerar solo un tipo de sensor, siendo en este caso conveniente permitir a los sensores que requieren cadena de medida acceder a la mayor cantidad de periféricos analógicos disponibles del \gls{mcu} (figura~\ref{fig:con_main}).

\begin{figure}[H]
    \centering
    \caption{Conector base propuesto para el nodo}
    \label{fig:con_main}
    \includegraphics[width=0.25\textwidth]{conector_principal}

    \nota{Los seis pines inferiores son la versión reducida; los demás hacen parte de la versión extendida.}

    \vspace{.3\baselineskip}
    \propia
\end{figure}

Dado que ambos conectores presentan pines en común, se puede escoger cualquiera de los dos según la complejidad del circuito esclavo por controlar.

\begin{table}[H]
\centering
\caption{Disposición sugerida de las funcionalidades de los pines}
\footnotesize
\label{tab:conector_pines}
\begin{tabular}{llllll}
\toprule
\textbf{Pin} & \textbf{Tipo} & \textbf{Función por defecto} & \textbf{Pin} & \textbf{Tipo} & \textbf{Función por defecto}          \\ \midrule
1   & PWR  & VCC, 2.8 V - 3.6 V  & 2   & PWR  & GND                                                                      \\ \midrule
3   & I/O  & VREFH                                                           & 4   & I/O  & VREFL                                                                    \\ \midrule
5   & I/O  & --Bloqueado--                                                   & 6   & I/O  & \parbox{3cm}{(CLK*/SCK)Op2,\\salida}        \\ \midrule
7   & I/O  & Op1, salida                                                     & 8   & I/O  & \parbox{3cm}{(MOSI*/Tx*)Op2,\\inversora}    \\ \midrule
9   & I/O  & Op1, inversora      & 10  & I/O  & \parbox{3cm}{(MISO*/Rx*)Op2,\\no inversora} \\ \midrule
11  & I/O  & Op1, no inversora   & 12  & I/O  & SS*(SDT)                                                                 \\ \midrule
13  & I/O  & Tx*                                                             & 14  & I/O  & MOSI*                                                                    \\ \midrule
15  & I/O  & Rx*                                                             & 16  & I/O  & MISO*                                                                    \\ \midrule
17  & I/O  & SDT*                                                            & 18  & I/O  & SCK/CLK*                                                                 \\ \bottomrule
\end{tabular}

\pieimagensimple{* Pines para periféricos digitales, que pueden ser asignados a otros pines, y entre paréntesis para el conector base.}
\end{table}

%%%%%%%%%%%%%%%%%%%%%%%%%%%%%%%%%%%%%%%%%%%
\subsection[Interfaz de \emph{software}: configuración mediante ID de núcleo y placa]{Interfaz de \emph{software}: configuración\\ mediante ID de núcleo y placa}

En aras de que el usuario final realice la mayor cantidad de configuración desde una interfaz gráfica y con un mínimo esfuerzo, se plantea una aplicación en que se ingresan dos identificadores: uno del núcleo y otro del circuito de aplicación.

\subsubsection{Núcleo}

Cada núcleo, conformado por un microcontrolador y un módem LoRaWAN, estará asociado a un identificador universal del núcleo (CoreID). Este identificador permite su indexación en una base de datos, en la que se almacenan, entre otros, los siguientes elementos particulares:
\begin{itemize}
    \item Microcontrolador principal.
    \item Marca y modelo del módem LoRaWAN.
    \item Tipo de antena.
\end{itemize}

\subsubsection{Circuito de aplicación}

Cada posible configuración eléctrica tendrá asignado un identificador de catálogo (CatID), repetible solamente entre el mismo tipo de circuito. De tal manera, cada cambio que altere la funcionalidad del \emph{hardware}, modificando el circuito, requerirá obligatoriamente un nuevo CatID. Entre los cambios notables, se presentan los siguientes:
\begin{itemize}
    \item Agregar/retirar componentes.
    \item Cambiar valores de componentes alterando las impedancias.
    \item Cambiar las posiciones relativas de los conectores.
\end{itemize}

Cambios en el encaminamiento de pistas o en la orientación de componentes discretos que no afecten el funcionamiento de los dispositivos generalmente no implicarán la creación de un nuevo CatID.

%%%%%%%%%%%%%%%%%%%%%%%%%%%%%%%%%%%%%%%%%%%
\section{Abstracción de interfaces eléctricas}

En el mercado existe gran variedad de sensores, cuyas variables de medida pueden reflejarse análogamente en propiedades eléctricas como algún voltaje, corriente, impedancia, o alguna combinación en particular. Del mismo modo, estos sensores, que incluyen gran parte de la cadena de medida y que presentan la variable mediante una interfaz digital, tienen su propio orden para la carga útil de cada una de estas interfaces.

Se pretende que el procesamiento de la información en el nodo sea mínimo, con el fin de aprovechar al máximo la energía disponible, sobre todo al venir de baterías o fuentes no disponibles la mayor parte del tiempo, trasladando el procesamiento intensivo a quienes hagan uso de la aplicación a la cual se enviarán los datos de los sensores, casi en crudo.
